\documentclass[
  man,
  longtable,
  nolmodern,
  notxfonts,
  notimes,
  colorlinks=true,linkcolor=blue,citecolor=blue,urlcolor=blue]{apa7}

\usepackage{amsmath}
\usepackage{amssymb}




\RequirePackage{longtable}
\RequirePackage{threeparttablex}

\makeatletter
\renewcommand{\paragraph}{\@startsection{paragraph}{4}{\parindent}%
	{0\baselineskip \@plus 0.2ex \@minus 0.2ex}%
	{-.5em}%
	{\normalfont\normalsize\bfseries\typesectitle}}

\renewcommand{\subparagraph}[1]{\@startsection{subparagraph}{5}{0.5em}%
	{0\baselineskip \@plus 0.2ex \@minus 0.2ex}%
	{-\z@\relax}%
	{\normalfont\normalsize\bfseries\itshape\hspace{\parindent}{#1}\textit{\addperi}}{\relax}}
\makeatother




\usepackage{longtable, booktabs, multirow, multicol, colortbl, hhline, caption, array, float, xpatch}
\usepackage{subcaption}


\renewcommand\thesubfigure{\Alph{subfigure}}
\setcounter{topnumber}{2}
\setcounter{bottomnumber}{2}
\setcounter{totalnumber}{4}
\renewcommand{\topfraction}{0.85}
\renewcommand{\bottomfraction}{0.85}
\renewcommand{\textfraction}{0.15}
\renewcommand{\floatpagefraction}{0.7}

\usepackage{tcolorbox}
\tcbuselibrary{listings,theorems, breakable, skins}
\usepackage{fontawesome5}

\definecolor{quarto-callout-color}{HTML}{909090}
\definecolor{quarto-callout-note-color}{HTML}{0758E5}
\definecolor{quarto-callout-important-color}{HTML}{CC1914}
\definecolor{quarto-callout-warning-color}{HTML}{EB9113}
\definecolor{quarto-callout-tip-color}{HTML}{00A047}
\definecolor{quarto-callout-caution-color}{HTML}{FC5300}
\definecolor{quarto-callout-color-frame}{HTML}{ACACAC}
\definecolor{quarto-callout-note-color-frame}{HTML}{4582EC}
\definecolor{quarto-callout-important-color-frame}{HTML}{D9534F}
\definecolor{quarto-callout-warning-color-frame}{HTML}{F0AD4E}
\definecolor{quarto-callout-tip-color-frame}{HTML}{02B875}
\definecolor{quarto-callout-caution-color-frame}{HTML}{FD7E14}

%\newlength\Oldarrayrulewidth
%\newlength\Oldtabcolsep


\usepackage{hyperref}




\providecommand{\tightlist}{%
  \setlength{\itemsep}{0pt}\setlength{\parskip}{0pt}}
\usepackage{longtable,booktabs,array}
\usepackage{calc} % for calculating minipage widths
% Correct order of tables after \paragraph or \subparagraph
\usepackage{etoolbox}
\makeatletter
\patchcmd\longtable{\par}{\if@noskipsec\mbox{}\fi\par}{}{}
\makeatother
% Allow footnotes in longtable head/foot
\IfFileExists{footnotehyper.sty}{\usepackage{footnotehyper}}{\usepackage{footnote}}
\makesavenoteenv{longtable}

\usepackage{graphicx}
\makeatletter
\newsavebox\pandoc@box
\newcommand*\pandocbounded[1]{% scales image to fit in text height/width
  \sbox\pandoc@box{#1}%
  \Gscale@div\@tempa{\textheight}{\dimexpr\ht\pandoc@box+\dp\pandoc@box\relax}%
  \Gscale@div\@tempb{\linewidth}{\wd\pandoc@box}%
  \ifdim\@tempb\p@<\@tempa\p@\let\@tempa\@tempb\fi% select the smaller of both
  \ifdim\@tempa\p@<\p@\scalebox{\@tempa}{\usebox\pandoc@box}%
  \else\usebox{\pandoc@box}%
  \fi%
}
% Set default figure placement to htbp
\def\fps@figure{htbp}
\makeatother


% definitions for citeproc citations
\NewDocumentCommand\citeproctext{}{}
\NewDocumentCommand\citeproc{mm}{%
  \begingroup\def\citeproctext{#2}\cite{#1}\endgroup}
\makeatletter
 % allow citations to break across lines
 \let\@cite@ofmt\@firstofone
 % avoid brackets around text for \cite:
 \def\@biblabel#1{}
 \def\@cite#1#2{{#1\if@tempswa , #2\fi}}
\makeatother
\newlength{\cslhangindent}
\setlength{\cslhangindent}{1.5em}
\newlength{\csllabelwidth}
\setlength{\csllabelwidth}{3em}
\newenvironment{CSLReferences}[2] % #1 hanging-indent, #2 entry-spacing
 {\begin{list}{}{%
  \setlength{\itemindent}{0pt}
  \setlength{\leftmargin}{0pt}
  \setlength{\parsep}{0pt}
  % turn on hanging indent if param 1 is 1
  \ifodd #1
   \setlength{\leftmargin}{\cslhangindent}
   \setlength{\itemindent}{-1\cslhangindent}
  \fi
  % set entry spacing
  \setlength{\itemsep}{#2\baselineskip}}}
 {\end{list}}
\usepackage{calc}
\newcommand{\CSLBlock}[1]{\hfill\break\parbox[t]{\linewidth}{\strut\ignorespaces#1\strut}}
\newcommand{\CSLLeftMargin}[1]{\parbox[t]{\csllabelwidth}{\strut#1\strut}}
\newcommand{\CSLRightInline}[1]{\parbox[t]{\linewidth - \csllabelwidth}{\strut#1\strut}}
\newcommand{\CSLIndent}[1]{\hspace{\cslhangindent}#1}





\usepackage{newtx}

\defaultfontfeatures{Scale=MatchLowercase}
\defaultfontfeatures[\rmfamily]{Ligatures=TeX,Scale=1}





\title{Climate Scientists and mental health: Wellbeing, burnout and hope
for the future}


\shorttitle{Climate Scientists and mental health}


\usepackage{etoolbox}









\authorsnames[{1},{2,1},{1,3},{1},{1,4}]{Amy Lykins,Suzanne M.
Cosh,Nicola Paul,Jude Fox,Deborah Apthorp}







\authorsaffiliations{
{School of Psychology, University of New England},{School of
Psychology, University of Adelaide},{UNE Business School, University of
New England},{School of Medicine and Psychology, Australian National
University}}




\leftheader{Lykins, Cosh, Paul, Fox and Apthorp}



\abstract{Abstract goes here  \section{Impact Statement} \noindent This
study reveals important information about the wellbeing of climate
scientists }

\keywords{Climate change, Climate scientists, Climate anxiety}

\authornote{\par{\addORCIDlink{Amy
Lykins}{0000-0003-2930-3964}}\par{\addORCIDlink{Suzanne M.
Cosh}{0000-0002-8003-3704}}\par{\addORCIDlink{Nicola
Paul}{0000-0001-8423-8742}}\par{\addORCIDlink{Deborah
Apthorp}{0000-0001-5785-024X}} 
\par{ }
\par{This study was preregistered at the
\href{https://osf.io/xegr7/registrations}{OSF} Raw data and code for the
study are available on the OSF {[}link{]}      }
\par{Correspondence concerning this article should be addressed to Amy
Lykins, Email: \href{mailto:alykins@une.edu.au}{alykins@une.edu.au}}
}

\makeatletter
\let\endoldlt\endlongtable
\def\endlongtable{
\hline
\endoldlt
}
\makeatother
\RequirePackage{longtable}
\DeclareDelayedFloatFlavor{longtable}{table}

\urlstyle{same}



\usepackage{booktabs}
\usepackage{longtable}
\usepackage{array}
\usepackage{multirow}
\usepackage{wrapfig}
\usepackage{float}
\usepackage{colortbl}
\usepackage{pdflscape}
\usepackage{tabu}
\usepackage{threeparttable}
\usepackage{threeparttablex}
\usepackage[normalem]{ulem}
\usepackage{makecell}
\usepackage{xcolor}
\makeatletter
\@ifpackageloaded{caption}{}{\usepackage{caption}}
\AtBeginDocument{%
\ifdefined\contentsname
  \renewcommand*\contentsname{Table of contents}
\else
  \newcommand\contentsname{Table of contents}
\fi
\ifdefined\listfigurename
  \renewcommand*\listfigurename{List of Figures}
\else
  \newcommand\listfigurename{List of Figures}
\fi
\ifdefined\listtablename
  \renewcommand*\listtablename{List of Tables}
\else
  \newcommand\listtablename{List of Tables}
\fi
\ifdefined\figurename
  \renewcommand*\figurename{Figure}
\else
  \newcommand\figurename{Figure}
\fi
\ifdefined\tablename
  \renewcommand*\tablename{Table}
\else
  \newcommand\tablename{Table}
\fi
}
\@ifpackageloaded{float}{}{\usepackage{float}}
\floatstyle{ruled}
\@ifundefined{c@chapter}{\newfloat{codelisting}{h}{lop}}{\newfloat{codelisting}{h}{lop}[chapter]}
\floatname{codelisting}{Listing}
\newcommand*\listoflistings{\listof{codelisting}{List of Listings}}
\makeatother
\makeatletter
\makeatother
\makeatletter
\@ifpackageloaded{caption}{}{\usepackage{caption}}
\@ifpackageloaded{subcaption}{}{\usepackage{subcaption}}
\makeatother

% From https://tex.stackexchange.com/a/645996/211326
%%% apa7 doesn't want to add appendix section titles in the toc
%%% let's make it do it
\makeatletter
\xpatchcmd{\appendix}
  {\par}
  {\addcontentsline{toc}{section}{\@currentlabelname}\par}
  {}{}
\makeatother

%% Disable longtable counter
%% https://tex.stackexchange.com/a/248395/211326

\usepackage{etoolbox}

\makeatletter
\patchcmd{\LT@caption}
  {\bgroup}
  {\bgroup\global\LTpatch@captiontrue}
  {}{}
\patchcmd{\longtable}
  {\par}
  {\par\global\LTpatch@captionfalse}
  {}{}
\apptocmd{\endlongtable}
  {\ifLTpatch@caption\else\addtocounter{table}{-1}\fi}
  {}{}
\newif\ifLTpatch@caption
\makeatother

\begin{document}

\maketitle




\setlength\LTleft{0pt}




\textsubscript{Source:
\href{https://deborahapthorp.github.io/Amy_ClimateScientists/index.qmd.html}{Article
Notebook}}

\textsubscript{Source:
\href{https://deborahapthorp.github.io/Amy_ClimateScientists/index.qmd.html}{Article
Notebook}}

\textsubscript{Source:
\href{https://deborahapthorp.github.io/Amy_ClimateScientists/index.qmd.html}{Article
Notebook}}

\textsubscript{Source:
\href{https://deborahapthorp.github.io/Amy_ClimateScientists/index.qmd.html}{Article
Notebook}}

\textsubscript{Source:
\href{https://deborahapthorp.github.io/Amy_ClimateScientists/index.qmd.html}{Article
Notebook}}

\textsubscript{Source:
\href{https://deborahapthorp.github.io/Amy_ClimateScientists/index.qmd.html}{Article
Notebook}}

\textsubscript{Source:
\href{https://deborahapthorp.github.io/Amy_ClimateScientists/index.qmd.html}{Article
Notebook}}

\textsubscript{Source:
\href{https://deborahapthorp.github.io/Amy_ClimateScientists/index.qmd.html}{Article
Notebook}}

The threat of climate change's consequences to human physical health
have long been recognised, whilst the threat to mental health has
received increasing attention over the past 10-15 years
(\citeproc{ref-doherty_psychological_2011}{Doherty \& Clayton, 2011}).
Researchers have proposed these effects to occur via one or more of
three pathways: 1) direct exposure to extreme weather events (e.g.,
wildfires, floods); 2) indirect effects of extreme weather and long-term
climatic changes (e.g., social unrest, forced migration); and 3)
distress associated with the awareness of climate change and its
anticipated consequences (\citeproc{ref-berry_climate_2010}{Berry et
al., 2010}), often referred to as `climate anxiety' or `eco-anxiety'
(\citeproc{ref-clayton_climate_2020}{Clayton, 2020}). Though there is no
generally agreed-upon definition of this construct
(\citeproc{ref-van_dijk_limited_2025}{Dijk et al., 2025}), research has
flourished in this space in recent years, showing reliably consistent
relationships with symptoms of depression, anxiety, and stress
(\citeproc{ref-cosh_relationship_2024}{Cosh et al., 2024}), as well as
{[}MORE HERE{]}

Some groups (e.g., children and young people) are considered to be at
disproportionately high risk of experiencing psychological distress
related to climate change in anticipation of futures that feel uncertain
and dangerous (\citeproc{ref-crandon_socialecological_2022}{Crandon et
al., 2022}). One other group proposed to be at high risk of
climate-related distress comprises people whose work brings them into
regular contact with the effects of climate change: namely, climate
scientists (\citeproc{ref-pihkala_cost_2020}{Pihkala, 2020}).

Despite potential interest, little empirical data currently exist on the
mental health and well-being of the world's climate scientists, and that
which does exist largely relies on relatively small sample sizes.
Several studies analysed 44 letters written by climate scientists posted
on Duggan's ``Is this how you feel?'' website
(\url{www.isthishowyoufeel.com}), reporting frustration and sadness as
most frequently mentioned related to climate change
(\citeproc{ref-clayton_mental_2018}{Clayton, 2018}), along with the
strain associated with managing multiple identities (e.g., personal,
parental, scientist, messenger) in the context of their work and lives
(\citeproc{ref-clayton_mental_2018}{Clayton, 2018}). Similar themes were
found in two qualitative studies. Head and Harada
(\citeproc{ref-head_keeping_2017}{2017}) interviewed 13 Australian
climate scientists, observing the tension between the need to remain
dispassionate and objective in their work, yet on a personal level
continuing to experience painful emotions associated with aspects of
their work (e.g., evidence of advancing climate change, receiving hate
mail and death threats), often resulting in the need to compartmentalise
their work from the rest of their lives. Renouf
(\citeproc{ref-renouf_making_2021}{2021}) {[}'s{]} interviews with 16
climate experts across 12 countries also revealed a range of negative
emotions experienced, including anxiety, sadness, and rage.
Compartmentalising efforts were less successful in this sample, with
findings showing climate change to have permeated these experts'
personal lives significantly.

Critically, research has also revealed positive emotions experienced by
climate scientists, with reports of hope, optimism, fascination, and
determination (\citeproc{ref-clayton_mental_2018}{Clayton, 2018}), along
with enjoyment found in the work (\citeproc{ref-head_keeping_2017}{Head
\& Harada, 2017}). Unfortunately, the balance of reports skews negative
(\citeproc{ref-haddaway_safe_2023}{Haddaway \& Duggan, 2023}), with
positive emotions sometimes only being reported when asked about
specifically (\citeproc{ref-renouf_making_2021}{Renouf, 2021}). A common
theme amongst research in this area is increasing pessimism that the
required changes will be made in time to avoid the most severe impacts
of climate change (\citeproc{ref-head_keeping_2017}{Head \& Harada,
2017}; \citeproc{ref-renouf_making_2021}{Renouf, 2021};
\citeproc{ref-tollefson_top_2021}{Tollefson, 2021}).

Early efforts into understanding how climate scientists cope with these
feelings and the different roles and responsibilities they have in life
suggested an emphasis on dispassion, denial, or suppression of these
emotions, as well as avoiding consideration of the implications of their
findings (\citeproc{ref-head_keeping_2017}{Head \& Harada, 2017}). More
recently, a study of 215 Lithuanian climate experts revealed both
problem-focused (i.e., direct efforts in solving climate change issues)
and emotion-focused (i.e., delving into feelings, seeking external
support) coping strategies preferred over avoidance-based coping (i.e.,
trying not to think about it;
\citeproc{ref-jovarauskaite_emotional_2021}{Jovarauskaite \& Böhm,
2021}). Though these results are yet to be replicated, they provide a
unique and testable assessment of how climate scientists may be coping
with feelings of anxiety, anger, and pessimism.

A recent scoping review emphasised the emotional burden climate
scientists may experience given their frequent and direct proximity with
the newest and most threatening information about climate change
(\citeproc{ref-calabria_scoping_2024}{Calabria \& Marks, 2024}).
Combined with a work environment that encourages dispassionate
objectivity, alongside often regular and vitriolic personal and
professional attacks and a worldwide relative lack of action to curb
climate change, climate scientists are uniquely at risk of moral
distress and moral injury (\citeproc{ref-calabria_scoping_2024}{Calabria
\& Marks, 2024}). These feelings have been observed in young people
(\citeproc{ref-henritze_moral_2023}{Henritze et al., 2023}); however, in
this population, this has been associated with a perceived lack of
agency to effect change, which may not be as much a factor for climate
scientists who are actively working on the problem
(\citeproc{ref-calabria_scoping_2024}{Calabria \& Marks, 2024}). Given
that higher perceived efficacy in mitigating the effects of climate
change may support mental health and well-being in the face of climate
distress (\citeproc{ref-ojala_anxiety_2021}{Ojala et al., 2021}), it is
also possible that engaging in climate science itself serves to protect
climate scientists' well-being in the face of these challenges. However,
to our knowledge, no one has yet examined these relationships in this
population.

To address this knowledge gap, we aimed to recruit a large,
international sample of climate scientists to investigate experiences of
personal and work-related burnout, psychological distress, and
well-being in people engaged in some aspect of climate science. To
examine potential differences across time, we collected data via a
mixed-methods online survey in two waves eight years apart (2017, 2025).
In this paper, we focus only on the quantitative data. Our hypotheses
and analytic strategy were pre-registered on the
\href{https://osf.io/xegr7/registrations}{OSF}. Based on the extant
literature, we hypothesised:

\begin{enumerate}
\def\labelenumi{\arabic{enumi}.}
\tightlist
\item
  Mental ill-health, as assessed by the DASS-21, will be significantly
  poorer (i.e., greater endorsement of symptoms) in Time 2 than in Time
  1
\item
  Personal burnout will be significantly higher in Time 2 than in Time 1
\item
  US-based participants will have lower well-being than non-US-based
  scientists
\item
  Personal burnout will be significantly higher than work burnout
\item
  Problem-focused coping will be associated with greater well-being than
  emotion-focused coping or avoidance
\item
  Organisational support will moderate the relationship between work
  burnout and well-being.
\end{enumerate}

\section{Results}\label{results}

Hypothesis 1 was that mental ill-health, as assessed by the DASS-21,
would be significantly poorer (i.e., greater endorsement of symptoms) at
Time 2 than in Time 1. To test this, we carried out three non-parametric
t-tests (Mann-Whitney U) on the three subscales of the DASS-21,
Depression, Anxiety, and Stress, since the data were heavily
right-skewed. This comparison was significant only for anxiety,
\(W = 10,046.00\), \(p = .005\), but not for depression,
\(W = 11,153.50\), \(p = .075\) or stress, \(W = 11,806.50\),
\(p = .217\). Thus Hypothesis 1 was partially supported. These
differences are illustrated in Figure~\ref{fig-DASS}.

\phantomsection\label{cell-fig-DASS}
\begin{figure}[H]

{\caption{{Mental health scores at Time 1 and Time 2 for the three
subscales of the DASS-21}{\label{fig-DASS}}}}

\pandocbounded{\includegraphics[keepaspectratio]{index_files/figure-pdf/fig-DASS-1.pdf}}

\noindent \emph{Note.} This~is a raincloud plot showing individual
scores and distributions for both time points.

\end{figure}

\textsubscript{Source:
\href{https://deborahapthorp.github.io/Amy_ClimateScientists/index.qmd.html}{Article
Notebook}}

Hypothesis 2 was that personal burnout at Time 2 would be higher than at
Time 1. Since the data for these scales were more normally distributed,
we tested this with a Welch's single-tailed independent samples t-test,
since the hypothesis was directional and the variances were unequal. We
also tested work burnout, using a two-tailed test since this hypothesis
was not preregistered. Personal burnout as measured by the Copenhagen
Burnout Inventory was significantly higher at Time 2 compared to Time 1,
\(\Delta M = -7.57\), 95\% CI \([-\infty, -3.49]\),
\(t(87.26) = -3.09\), \(p = .001\). Work burnout was also significantly
higher, \(\Delta M = -8.93\), 95\% CI \([-13.62, -4.24]\),
\(t(88.99) = -3.78\), \(p < .001\). This difference is illustrated in
Figure~\ref{fig-CBI}.

\phantomsection\label{cell-fig-CBI}
\begin{figure}[H]

{\caption{{Burnout scores for personal and work subscales of the
Copenhagen Burnout Inventory at Times 1 and 2.}{\label{fig-CBI}}}}

\pandocbounded{\includegraphics[keepaspectratio]{index_files/figure-pdf/fig-CBI-1.pdf}}

\noindent \emph{Note.} This~is a raincloud plot showing individual
scores and distributions for both time points.

\end{figure}

\textsubscript{Source:
\href{https://deborahapthorp.github.io/Amy_ClimateScientists/index.qmd.html}{Article
Notebook}}

Hypothesis 3 was that, at Time 2, wellbeing would be significantly lower
for US-based participants compared to their non-US counterparts. To test
this, we first grouped the data into those who responded that they were
US-based for work, using string-matching to ``US'', ``USA'' or ``United
States'', and those who were based elsewhere. We then compared these
groups on the wellbeing scale (SWEMS) using a general linear model,
controlling for age and gender. The results are presented in
Figure~\ref{fig-USA} and Table~\ref{tbl-USA}. Hypothesis 3 was not
supported; the difference in wellbeing was not significant,
\(b = -0.11\), 95\% CI \([-0.23, 0.02]\), \(t = -1.61\), \(p = .108\).

\phantomsection\label{cell-fig-USA}
\begin{figure}[H]

{\caption{{Wellbeing scores for participants who reported working in the
USA compared to those from other nations.}{\label{fig-USA}}}}

\pandocbounded{\includegraphics[keepaspectratio]{index_files/figure-pdf/fig-USA-1.pdf}}

\noindent \emph{Note.} This~is a raincloud plot showing individual
scores and distributions.

\end{figure}

\textsubscript{Source:
\href{https://deborahapthorp.github.io/Amy_ClimateScientists/index.qmd.html}{Article
Notebook}}

\begin{table}

{\caption{{}{\label{tbl-USA}}}
\vspace{-20pt}}

\begin{center}
\begin{threeparttable}

\caption{\label{tab:tbl-USA}Regression table for the GLM of wellbeing scores for participants within and outside the USA, controlling for age and gender}

\begin{tabular}{lllll}
\toprule
Predictor & \multicolumn{1}{c}{$b$} & \multicolumn{1}{c}{95\% CI} & \multicolumn{1}{c}{$t$} & \multicolumn{1}{c}{$p$}\\
\midrule
Intercept & 2.66 & {}[2.38, 2.95] & 18.57 & < .001\\
USATRUE & -0.11 & {}[-0.23, 0.02] & -1.61 & .108\\
Age & 0.01 & {}[0.01, 0.02] & 4.38 & < .001\\
Gender & 0.14 & {}[0.04, 0.24] & 2.74 & .006\\
\bottomrule
\end{tabular}

\end{threeparttable}
\end{center}

\end{table}

\textsubscript{Source:
\href{https://deborahapthorp.github.io/Amy_ClimateScientists/index.qmd.html}{Article
Notebook}}

Hypothesis 4 was that personal burnout would be higher than work burnout
across both waves. This analysis was a one-tailed paired-samples t-test,
since all participants completed both subscales of the Copenhagen
Burnout Inventory, and our hypothesis was directional. This hypothesis
was supported, with a mean difference of 3.31 percentage points,
\(M_D = 3.31\), 95\% CI \([2.38, \infty]\), \(t(485) = 5.90\),
\(p < .001\). The results are illustrated in
Figure~\ref{fig-CBI-paired}.

\phantomsection\label{cell-fig-CBI-paired}
\begin{figure}[H]

{\caption{{Personal compared to work burnout (measured as a percentage
from 0 to 100) on the Copenhagen Burnout Inventory, for participants
across both waves.}{\label{fig-CBI-paired}}}}

\pandocbounded{\includegraphics[keepaspectratio]{index_files/figure-pdf/fig-CBI-paired-1.pdf}}

\vspace{-20pt}
\noindent \emph{Note.} This~is a raincloud plot showing individual
scores and distributions, with grey lines connecting eqch participant's
data for the two measures.

\end{figure}

\textsubscript{Source:
\href{https://deborahapthorp.github.io/Amy_ClimateScientists/index.qmd.html}{Article
Notebook}}

Hypothesis 5 was that problem-focused coping as a strategy would be
associated with greater well-being than emotion-focused coping or
avoidance. For this analysis we conducted a multiple regression with
psychological wellbeing (as measured by the SWEMBS) as the outcome
variable, with the three coping subscales as predictors, and age and
gender as control variables. Hypothesis 5 was not supported; after
controlling for age and gender, only avoidance was associated (strongly,
and negatively) with wellbeing, \(b = -0.15\), 95\% CI
\([-0.23, -0.07]\), \(t(400) = -3.86\), \(p < .001\). The whole model
was significant, \(R^2 = .12\), 90\% CI \([0.07, 0.17]\),
\(F(5, 400) = 11.42\), \(p < .001\). The full results are presented in
Table~\ref{tbl-hyp5}.

\begin{table}

{\caption{{}{\label{tbl-hyp5}}}
\vspace{-20pt}}

\begin{center}
\begin{threeparttable}

\caption{\label{tab:tbl-hyp5}Coefficients for the multiple regression model of coping styles on wellbeing as measured by the SWEMWS scale}

\begin{tabular}{llllll}
\toprule
Predictor & \multicolumn{1}{c}{$b$} & \multicolumn{1}{c}{95\% CI} & \multicolumn{1}{c}{$t$} & \multicolumn{1}{c}{$\mathit{df}$} & \multicolumn{1}{c}{$p$}\\
\midrule
Intercept & -0.01 & {}[-0.07, 0.05] & -0.36 & 400 & .721\\
Coping PF score & 0.03 & {}[-0.03, 0.10] & 1.06 & 400 & .288\\
Coping EF score & 0.04 & {}[-0.01, 0.10] & 1.46 & 400 & .144\\
Coping avoidance score & -0.15 & {}[-0.23, -0.07] & -3.86 & 400 & < .001\\
Age & 0.01 & {}[0.00, 0.01] & 2.87 & 400 & .004\\
Gender & 0.13 & {}[0.03, 0.22] & 2.51 & 400 & .012\\
\bottomrule
\end{tabular}

\end{threeparttable}
\end{center}

\end{table}

\textsubscript{Source:
\href{https://deborahapthorp.github.io/Amy_ClimateScientists/index.qmd.html}{Article
Notebook}}

For Hypothesis 6, that organisational support would moderate the
relationship between work burnout and wellbeing, we conducted a multiple
regression analysis with psychological wellbeing as the DV, work burnout
and perceived organisational support (SPOS score) as the IVs, and an
interaction term between work burnout and SPOS, with age and gender as
control variables. From Table~\ref{tbl-hyp6}, it is clear that
Hypothesis 6 is also not supported; although work burnout had a large
and highly significant negative effect on wellbeing, \(b = -0.02\), 95\%
CI \([-0.02, -0.01]\), \(t(400) = -11.69\), \(p < .001\), and perceived
organisational support has a small marginal effect, \(b = 0.03\), 95\%
CI \([0.00, 0.06]\), \(t(400) = 2.08\), \(p = .038\), the interaction
between them was not significant, \(b = 0.00\), 95\% CI
\([0.00, 0.00]\), \(t(400) = -0.38\), \(p = .703\). The whole model was
significant, \(R^2 = .36\), 90\% CI \([0.29, 0.42]\),
\(F(5, 400) = 45.70\), \(p < .001\).

\begin{table}

{\caption{{}{\label{tbl-hyp6}}}
\vspace{-20pt}}

\begin{center}
\begin{threeparttable}

\caption{\label{tab:tbl-hyp6}Coefficients for the multiple regression model of work burnout and perceived organisational support on wellbeing as measured by the SWEMWS scale}

\begin{tabular}{llllll}
\toprule
Predictor & \multicolumn{1}{c}{$b$} & \multicolumn{1}{c}{95\% CI} & \multicolumn{1}{c}{$t$} & \multicolumn{1}{c}{$\mathit{df}$} & \multicolumn{1}{c}{$p$}\\
\midrule
Intercept & -0.01 & {}[-0.06, 0.04] & -0.35 & 400 & .727\\
Cbi work & -0.02 & {}[-0.02, -0.01] & -11.69 & 400 & < .001\\
SPOS score & 0.03 & {}[0.00, 0.06] & 2.08 & 400 & .038\\
Age & 0.00 & {}[0.00, 0.01] & 0.93 & 400 & .355\\
Gender & 0.05 & {}[-0.03, 0.13] & 1.16 & 400 & .248\\
Cbi work $\times$ SPOS score & 0.00 & {}[0.00, 0.00] & -0.38 & 400 & .703\\
\bottomrule
\end{tabular}

\end{threeparttable}
\end{center}

\end{table}

\textsubscript{Source:
\href{https://deborahapthorp.github.io/Amy_ClimateScientists/index.qmd.html}{Article
Notebook}}

\section{Discussion}\label{discussion}

\subsection{Limitations and Future
Directions}\label{limitations-and-future-directions}

\subsection{Conclusion}\label{conclusion}

\section{Online Methods}\label{online-methods}

\subsection{Participants}\label{participants}

Wave 1 participants were recruited between 1 March and 23 June, 2017. In
total, 61 participants (51 male, 10 female) with a mean age of 54.2
years (SD = 13.0, range: 23 -- 80 years) completed the online survey.
More than 90\% reported their highest relevant degree obtained as a PhD,
with 6.6\% reporting master's degree, and 3.3\% reporting other (e.g.,
Scientiae Doctor). Participants primarily resided in the United States
(n = 23), Australia (n = 15), and the United Kingdom (n = 13), with two
participants from Canada and one each from Brazil, Denmark, France,
Germany, Netherlands, New Zealand, Norway, and Switzerland. The majority
of participants reported conducting most of their research in their home
countries, but others mentioned `Europe wide,' along with France,
Greenland, and Republic of Korea. Participants reported between 3 and 55
years of experience in the field of climate science (M = 29.49 years, SD
= 12.34), and described their work as climate science (n = 18),
atmospheric science (n = 12), climate modelling (n = 5), oceanography or
marine science (n = 5), climate policy (n = 3), or with other terms (n =
16) including glaciology, ecosystems, and paleoclimatology. Wave 2
participants were recruited between 3 March and 1 July, 2025. Overall,
427 participants (250 male, 158 female, 19 `other' or prefer not to say)
with a mean age of 54.05 years (SD = 12.32, range = 27 -- 85) completed
the online survey. Highest degree obtained was Ph.D.~for 96\% of our
participants, followed by master's degree (2.3\%), bachelor's/honours
(0.4\%), and other (1.2\%; Doctor of Science, M.D., D.V.M.).

\subsection{Measures: Wave 1 and 2}\label{measures-wave-1-and-2}

\subsubsection{Demographics}\label{demographics}

Participants were asked a series of demographic questions including age,
gender identity, highest level of education attained, the country in
which they resided, the country/countries in which they conducted their
research, the type of organisation they worked for, and the nature of
their role within that organisation.

\subsubsection{Depression Anxiety Stress Scales
(DASS-21)}\label{depression-anxiety-stress-scales-dass-21}

The DASS-21 (\citeproc{ref-lovibond_manual_1995}{S. H. Lovibond \&
Lovibond, 1995}) is a 21-item self-report screening measure of
depression, anxiety, and stress symptoms experienced over the preceding
week. Item response options ranged from 0 (\emph{never}) to 3
(\emph{almost always}). The DASS-21 has well-validated clinical cut-offs
identified for mild, moderate, severe, and extremely severe
presentations of symptoms on each subscale, with moderate symptoms
established as a baseline for clinical relevance. The DASS-21
consistently has been shown to have good psychometric properties
(\citeproc{ref-lovibond_structure_1995}{P. F. Lovibond \& Lovibond,
1995}). Cronbach's alphas in the current samples were .89, .80, and .89
for depression, anxiety and stress respectively.

\subsubsection{Copenhagen Burnout Inventory
(CBI)}\label{copenhagen-burnout-inventory-cbi}

Two subscales of the CBI
(\citeproc{ref-kristensen_copenhagen_2005}{Kristensen et al., 2005})
were used to assess the personal burnout (e.g., \emph{How often do you
feel worn out?}; 6 items) and work-related burnout (e.g., \emph{Is your
work emotionally exhausting?}; 7 items) experiences of our participants.
Likert-scale response options ranged from (\emph{never/almost never} or
\emph{to a very low degree}) to 100 (\emph{always or to a very high
degree}), with three other response options available valued at 25, 50,
and 75. Kristensen et al.
(\citeproc{ref-kristensen_copenhagen_2005}{2005}) reported very high
internal reliability and occupational differentiation, indicating good
validity. Correlations with assessments of fatigue and psychological
well-being were found in the predicted directions, and burnout scores
predicted other self-reported well-being indicators such as future
sickness absence, sleep problems, and intention to quit. Cronbach's
alphas in the current samples for personal and work-related burnout were
.91 and .89, respectively.

\subsection{Measures: Wave 2}\label{measures-wave-2}

\subsubsection{Short Warwick-Edinburgh Mental Well-being Scale
(SWEMWBS).}\label{short-warwick-edinburgh-mental-well-being-scale-swemwbs.}

The SWEMWBS (\citeproc{ref-shah_short_2021}{Shah et al., 2021}) is a
widely used 7-item shortened version of the original Warwick-Edinburgh
Mental Well-being Scale
(\citeproc{ref-tennant_warwick-edinburgh_2007}{Tennant et al., 2007})
designed to assess positive experiences of subjective mental well-being
over the previous 2 weeks (e.g., \emph{I've been dealing with problems
well}). Likert-scale response options ranged from 1 (\emph{none of the
time}) to 5 (\emph{all of the time}). A recent systematic review
confirmed the SWEMWBS's good psychometric properties with respect to
validity and reliability (\citeproc{ref-perera_psychometric_2025}{Perera
et al., 2025}). Cronbach's alpha in the current sample was .84.

\subsubsection{Work and Meaning Inventory
(WAMI)}\label{work-and-meaning-inventory-wami}

Participants completed the 10-item WAMI
(\citeproc{ref-steger_measuring_2012}{Steger et al., 2012}) to assess
the degree to which they perceived their work as a subjectively
meaningful experience. The WAMI is comprised of three subscales:
positive meaning (e.g., \emph{I have a good sense of what makes my job
meaningful}; 4 items), meaning making through work (e.g., \emph{My work
helps me make sense of the world around me}; 3 items), and greater good
motivations (e.g., \emph{The work I do serves a greater purpose}; 3
items). Likert-scale response options ranged from 1 (\emph{absolutely
untrue}) to 5 (\emph{absolutely true}). Initial
(\citeproc{ref-steger_measuring_2012}{Steger et al., 2012}) and later
(\citeproc{ref-paola_evaluating_2023}{Paola et al., 2023}) assessments
of the WAMI indicate its strong psychometric properties with respect to
validity and reliability. Cronbach's alpha in the current sample was .89

\subsubsection{Survey of Perceived Organizational Support
(SPOS-8)}\label{survey-of-perceived-organizational-support-spos-8}

The SPOS-8 (\citeproc{ref-bardhoshi_psychometric_2023}{Bardhoshi et al.,
2023}) is an 8-item shortened version of Eisenberger et al.'s
(\citeproc{ref-eisenberger_perceived_1986}{1986}) 36-item measure which
was developed to assess an employee's general belief that their employer
is committed to him/her/them, values the employee's ongoing membership,
and is concerned about the employee's well-being (e.g., \emph{This
organisation really cares about my well-being}). Likert-scale response
options ranged from 0 (\emph{strongly disagree}) to 6 (\emph{strongly
agree}). Assessments of the longer version have shown unidimensionality,
reliability, and acceptable item-total correlations, and examinations of
the 8-item version have shown good reliability
(\citeproc{ref-bardhoshi_psychometric_2023}{Bardhoshi et al., 2023};
\citeproc{ref-hellman_reliability_2006}{Hellman et al., 2006}).
Cronbach's alpha in the current sample was .92.

\subsubsection{Hogg Eco-anxiety Scale
(HEAS)}\label{hogg-eco-anxiety-scale-heas}

The HEAS (\citeproc{ref-hogg_hogg_2021}{Hogg et al., 2021}) is a 13-item
scale designed to assess anxiety about a variety of environmental
conditions (e.g., climate change, environmental degradation, and
pollution) and one's personal negative impact on the planet across the
preceding two weeks. The HEAS is comprised of 4 subscales: affective
symptoms (e.g., \emph{not being able to stop or control worrying}; 4
items), rumination (e.g., \emph{unable to stop thinking about future
climate change and other global environmental problems}; 3 items),
behavioural symptoms (e.g., \emph{difficulty enjoying social situations
with family and friends}; 3 items) , and anxiety about one's negative
impact on the planet (e.g., \emph{feeling anxious that your personal
behaviours will do little to help fix the problem}; 3 items).
Likert-scale response options range from 0 (\emph{not at all}) to 3
(\emph{nearly every day}). The HEAS shows strong psychometric properties
including good test-retest reliability, along with differentiation from
more general stress, anxiety, and depression
(\citeproc{ref-hogg_validation_2024}{Hogg et al., 2024}). Cronbach's
alphas from the current sample were .89 (affective), .88 (rumination),
.83 (behavioural) and .89 (personal impact).

\subsubsection{Coping Strategies.}\label{coping-strategies.}

To assess how participants any distress they may experience related to
climate change, participants completed the 11-item coping strategies
scale for climate experts created by Jovarauskaite and Böhm
(\citeproc{ref-jovarauskaite_emotional_2021}{2021}). The scale breaks
down coping strategies into three types: problem-focused strategies
(e.g., \emph{{[}I{]} concentrate on the ways in which climate change can
be solved}; 3 items), emotion-focused strategies (e.g., \emph{{[}I{]}
tell others how I feel}; 4 items), and avoidance (e.g., \emph{{[}I{]}
pretend that climate change is not happening}; 4 items). Likert-scale
response options ranged from 1 (\emph{strongly disagree}) to 6
(\emph{strongly agree}). This scale was validated on a sample of
Lithuanian climate experts, with the confirmatory factor analysis
showing good model fit
(\citeproc{ref-jovarauskaite_emotional_2021}{Jovarauskaite \& Böhm,
2021}). Because this scale is relatively new, we also conducted a
confirmatory factor analysis to confirm its factor structure, with
results indicating an acceptable 3-factor fit once item 8 from the
avoidance subscale was removed (see Supplementary Materials). Cronbach's
alphas in our sample were .82 for problem-focused coping, .86 for
emotion-focused coping, and .71 for avoidance.

\subsection{Procedure}\label{procedure}

All participants were recruited for an international study of climate
scientists. Each participant was sent a personalised email invitation
with a unique link to the online survey, hosted by the survey company
Qualtrics. After providing informed consent, participants completed the
demographic questions, followed by the other study measures in a
randomised order. Participation took approximately 10 minutes.

In Wave 1, a member of the research team (J. F.) compiled a list of
names and email addresses of researchers who had published in the area
of climate science, including contributors to the Intergovernmental
Panel on Climate Change (IPCC) publications. Personalised invitations
were emailed to 538 individuals, with approximately 15\% of them
bouncing for various reasons (e.g., email address no longer being in
use, individual on leave). The research team also approached a number of
relevant governmental and non-governmental organisations, several of
whom agreed to share our invitation to their members. In Wave 2, a
different member of the research team (N. P.) conducted a National
Institutes of Health (NIH) Medline search in December, 2023 of journal
articles published with `climate change' as a MeSH (Medical Subject
Heading), resulting in more than 15,000 publications since 2009.
Separately, we created a list of the names of authors of the IPCC 6th
Assessment. After cross-checking and removing individuals listed more
than once, and finding the email addresses of IPCC authors that were not
listed in the Medline search via internet searches, our final contact
list comprised 7,446 names and email addresses of individuals who either
had published peer-reviewed journal articles in climate science or had
contributed to the latest IPCC assessment at the time of recruitment. We
first contacted all of these authors on 3 March, 2025 and sent one
reminder (5 June, 2025) to everyone on this list before closing the
survey in mid-July, 2025. We received approximately 500 email bounces on
each contact (i.e., 1000+ total). Both waves of data collection received
approval from the University of New England's Human Research Ethics
Committee prior to participant recruitment (W1: HE17-032; W2: HE23-011)

\subsection{Data Analysis}\label{data-analysis}

Data were analysed in R Statistical Software (Version 4.5.1) using R
Studio (version 2025.09.1+401). Analysis code is available on GitHub
(link). Analyses for each hypothesis were conducted according to the
preregistered plan: for Hypothesis 1 (\emph{Mental health at time 2 will
be poorer than at Time 1}), we used non-parametric t-tests (Mann-Whitney
U) for each of the DASS subscales to compare Wave 2 with Wave 1. For
Hypothesis 2 (\emph{Personal burnout at time 2 will be higher than at
Time 1}), since the assumption of normality for the residuals was met,
we used a one-tailed t-test, since the hypothesis was directional; we
also compared work burnout using a two-tailed t-test as an exploratory
analysis. For Hypothesis 3 (\emph{At Time 2, wellbeing will be
significantly lower for US-based participants compared to their non-US
counterparts}), we used a general linear model with psychological
well-being (SWEMWBS score) as the DV, and geographical location
(US-based vs.~non-US-based) as a factor, with age and gender as control
variables. For Hypothesis 4 (\emph{Personal burnout will be higher than
work burnout}), we used a paired-samples t-test on data from both time
points for all participants. For Hypothesis 5 (\emph{Problem-focused
coping will be associated with greater well-being than emotion-focused
coping or avoidance}) we used a multiple linear regression with
psychological well-being (SWEMWBS score) as the dependent variable, with
IVs of the three different coping styles (i.e., problem-focused,
emotion-focused, and avoidance), and age and gender as control
variables; this analysis only used the Wave 2 data, as these variables
were not collected in Wave 1. For Hypothesis 6 (\emph{Organisational
support will moderate the relationship between work burnout and
well-being}), we also used the Wave 2 data only, and we conducted a
multiple linear regression with psychological wellbeing (SWEMWBS score)
as the DV, organisational support (SPOS) and work burnout (CBI-work) as
predictors, and an interaction term between SPOS and CBI-work; all
variables were centred before running the model.

In addition to the preregistered analyses, we conducted a confirmatory
factor analysis on the coping scales
(\citeproc{ref-jovarauskaite_emotional_2021}{Jovarauskaite \& Böhm,
2021}), since it is a relatively new measure which has not been
extensively validated. For this, we used the R package lavaan
(\citeproc{ref-rosseel_lavaan_2025}{Rosseel et al., 2025}), since there
are established subscales onto which the items are designed to load
(problem-focused coping, emotion-focused coping and avoidance). We used
robust maximum likelihood for the estimator and full information maximum
likelihood for missing values.

\section{References}\label{references}

\phantomsection\label{refs}
\begin{CSLReferences}{1}{0}
\bibitem[\citeproctext]{ref-bardhoshi_psychometric_2023}
Bardhoshi, G., Erford, B. T., \& Um, B. (2023). Psychometric {Analysis}
of the {Survey} of {Perceived} {Organizational} {Support} ({SPOS}-8)
{Using} 1,005 {School} {Counselors}. \emph{Measurement and Evaluation in
Counseling and Development}, \emph{56}(3), 241--253.
\url{https://doi.org/10.1080/07481756.2022.2070074}

\bibitem[\citeproctext]{ref-berry_climate_2010}
Berry, H. L., Bowen, K., \& Kjellstrom, T. (2010). Climate change and
mental health: A causal pathways framework. \emph{International Journal
of Public Health}, \emph{55}(2), 123--132.
\url{https://doi.org/10.1007/s00038-009-0112-0}

\bibitem[\citeproctext]{ref-calabria_scoping_2024}
Calabria, L., \& Marks, E. (2024). A scoping review of the impact of
eco-distress and coping with distress on the mental health experiences
of climate scientists. \emph{Frontiers in Psychology}, \emph{15}.
\url{https://doi.org/10.3389/fpsyg.2024.1351428}

\bibitem[\citeproctext]{ref-clayton_mental_2018}
Clayton, S. (2018). Mental health risk and resilience among climate
scientists. \emph{Nature Climate Change}, \emph{8}(4), 260--261.
\url{https://doi.org/10.1038/s41558-018-0123-z}

\bibitem[\citeproctext]{ref-clayton_climate_2020}
Clayton, S. (2020). Climate anxiety: {Psychological} responses to
climate change. \emph{Journal of Anxiety Disorders}, \emph{74}, 102263.
\url{https://doi.org/10.1016/j.janxdis.2020.102263}

\bibitem[\citeproctext]{ref-cosh_relationship_2024}
Cosh, S. M., Ryan, R., Fallander, K., Robinson, K., Tognela, J., Tully,
P. J., \& Lykins, A. D. (2024). The relationship between climate change
and mental health: A systematic review of the association between
eco-anxiety, psychological distress, and symptoms of major affective
disorders. \emph{BMC Psychiatry}, \emph{24}(1), 833.
\url{https://doi.org/10.1186/s12888-024-06274-1}

\bibitem[\citeproctext]{ref-crandon_socialecological_2022}
Crandon, T. J., Scott, J. G., Charlson, F. J., \& Thomas, H. J. (2022).
A social--ecological perspective on climate anxiety in children and
adolescents. \emph{Nature Climate Change}, \emph{12}(2), 123--131.
\url{https://doi.org/10.1038/s41558-021-01251-y}

\bibitem[\citeproctext]{ref-van_dijk_limited_2025}
Dijk, S. van, Schie, K. van, Smeets, T., \& Mertens, G. (2025). Limited
consensus on what climate anxiety is: {Insights} from content overlap
analysis on 12 questionnaires. \emph{Journal of Anxiety Disorders},
\emph{109}, 102957. \url{https://doi.org/10.1016/j.janxdis.2024.102957}

\bibitem[\citeproctext]{ref-doherty_psychological_2011}
Doherty, T. J., \& Clayton, S. (2011). The psychological impacts of
global climate change. \emph{American Psychologist}, \emph{66}(4),
265--276. \url{https://doi.org/10.1037/a0023141}

\bibitem[\citeproctext]{ref-eisenberger_perceived_1986}
Eisenberger, R., Huntington, R., Hutchison, S., \& Sowa, D. (1986).
Perceived organizational support. \emph{Journal of Applied Psychology},
\emph{71}(3), 500--507. \url{https://doi.org/10.1037/0021-9010.71.3.500}

\bibitem[\citeproctext]{ref-haddaway_safe_2023}
Haddaway, N. R., \& Duggan, J. (2023). {``{Safe} {Spaces}''} and
{Community} {Building} for {Climate} {Scientists}, {Exploring}
{Emotions} {Through} a {Case} {Study}. \emph{Global Environmental
Psychology}, \emph{1}, 1--23. \url{https://doi.org/10.5964/gep.11347}

\bibitem[\citeproctext]{ref-head_keeping_2017}
Head, L., \& Harada, T. (2017). Keeping the heart a long way from the
brain: {The} emotional labour of climate scientists. \emph{Emotion,
Space and Society}, \emph{24}, 34--41.
\url{https://doi.org/10.1016/j.emospa.2017.07.005}

\bibitem[\citeproctext]{ref-hellman_reliability_2006}
Hellman, C. M., Fuqua, D. R., \& Worley, J. (2006). A {Reliability}
{Generalization} {Study} on the {Survey} of {Perceived} {Organizational}
{Support}: {The} {Effects} of {Mean} {Age} and {Number} of {Items} on
{Score} {Reliability}. \emph{Educational and Psychological Measurement},
\emph{66}(4), 631--642. \url{https://doi.org/10.1177/0013164406288158}

\bibitem[\citeproctext]{ref-henritze_moral_2023}
Henritze, E., Goldman, S., Simon, S., \& Brown, A. D. (2023). Moral
injury as an inclusive mental health framework for addressing climate
change distress and promoting justice-oriented care. \emph{The Lancet.
Planetary Health}, \emph{7}(3), e238--e241.
\url{https://doi.org/10.1016/S2542-5196(22)00335-7}

\bibitem[\citeproctext]{ref-hogg_validation_2024}
Hogg, T. L., Stanley, S. K., \& O'Brien, L. V. (2024). Validation of the
{Hogg} {Climate} {Anxiety} {Scale}. \emph{Climatic Change},
\emph{177}(6), 86. \url{https://doi.org/10.1007/s10584-024-03726-1}

\bibitem[\citeproctext]{ref-hogg_hogg_2021}
Hogg, T. L., Stanley, S. K., O'Brien, L. V., Wilson, M. S., \& Watsford,
C. R. (2021). The {Hogg} {Eco}-{Anxiety} {Scale}: {Development} and
validation of a multidimensional scale. \emph{Global Environmental
Change}, \emph{71}, 102391.
\url{https://doi.org/10.1016/j.gloenvcha.2021.102391}

\bibitem[\citeproctext]{ref-jovarauskaite_emotional_2021}
Jovarauskaite, L., \& Böhm, G. (2021). The emotional engagement of
climate experts is related to their climate change perceptions and
coping strategies. \emph{Journal of Risk Research}, \emph{24}(8),
941--957. \url{https://doi.org/10.1080/13669877.2020.1779785}

\bibitem[\citeproctext]{ref-kristensen_copenhagen_2005}
Kristensen, T. S., Borritz, M., Villadsen, E., \& Christensen, K. B.
(2005). The {Copenhagen} {Burnout} {Inventory}: {A} new tool for the
assessment of burnout. \emph{Work \& Stress}, \emph{19}(3), 192--207.
\url{https://doi.org/10.1080/02678370500297720}

\bibitem[\citeproctext]{ref-lovibond_structure_1995}
Lovibond, P. F., \& Lovibond, S. H. (1995). The structure of negative
emotional states: {Comparison} of the {Depression} {Anxiety} {Stress}
{Scales} ({DASS}) with the {Beck} {Depression} and {Anxiety}
{Inventories}. \emph{Behaviour Research and Therapy}, \emph{33}(3),
335--343. \url{https://doi.org/10.1016/0005-7967(94)00075-U}

\bibitem[\citeproctext]{ref-lovibond_manual_1995}
Lovibond, S. H., \& Lovibond, P. F. (1995). \emph{Manual for the
depression anxiety stress scales} (2nd ed). Psychology Foundation of
Australia.

\bibitem[\citeproctext]{ref-ojala_anxiety_2021}
Ojala, M., Cunsolo, A., Ogunbode, C. A., \& Middleton, J. (2021).
Anxiety, {Worry}, and {Grief} in a {Time} of {Environmental} and
{Climate} {Crisis}: {A} {Narrative} {Review}. \emph{Annual Review of
Environment and Resources}, \emph{46}(Volume 46, 2021), 35--58.
\url{https://doi.org/10.1146/annurev-environ-012220-022716}

\bibitem[\citeproctext]{ref-paola_evaluating_2023}
Paola, M., Rita, Z., \& Giuseppe, S. (2023). Evaluating meaningful work:
{Psychometric} properties of the {Work} and {Meaning} {Inventory}
({WAMI}) in {Italian} context. \emph{Current Psychology}, \emph{42}(15),
12756--12767. \url{https://doi.org/10.1007/s12144-021-02503-y}

\bibitem[\citeproctext]{ref-perera_psychometric_2025}
Perera, B. P. R., Wickremasinghe, A. R., \& De Za, T. a. P. (2025).
Psychometric properties of the {Warwick} {Edinburgh} {Mental}
{Well}-being {Scale}: A systematic review. \emph{Systematic Reviews},
\emph{14}(1), 149. \url{https://doi.org/10.1186/s13643-025-02897-x}

\bibitem[\citeproctext]{ref-pihkala_cost_2020}
Pihkala, P. (2020). The {Cost} of {Bearing} {Witness} to the
{Environmental} {Crisis}: {Vicarious} {Traumatization} and {Dealing}
with {Secondary} {Traumatic} {Stress} among {Environmental}
{Researchers}. \emph{Social Epistemology}, \emph{34}(1), 86--100.
\url{https://doi.org/10.1080/02691728.2019.1681560}

\bibitem[\citeproctext]{ref-renouf_making_2021}
Renouf, J. S. (2021). Making sense of climate change---the lived
experience of experts. \emph{Climatic Change}, \emph{164}(1), 14.
\url{https://doi.org/10.1007/s10584-021-02986-5}

\bibitem[\citeproctext]{ref-rosseel_lavaan_2025}
Rosseel, Y., Jorgensen, T. D., Wilde, L. D., Oberski, D., Byrnes, J.,
Vanbrabant, L., Savalei, V., Merkle, E., Hallquist, M., Rhemtulla, M.,
Katsikatsou, M., Barendse, M., Rockwood, N., Scharf, F., Du, H., Jamil,
H., \& Classe, F. (2025). \emph{Lavaan: {Latent} {Variable} {Analysis}}.
\url{https://cran.r-project.org/web/packages/lavaan/index.html}

\bibitem[\citeproctext]{ref-shah_short_2021}
Shah, N., Cader, M., Andrews, B., McCabe, R., \& Stewart-Brown, S. L.
(2021). Short {Warwick}-{Edinburgh} {Mental} {Well}-being {Scale}
({SWEMWBS}): Performance in a clinical sample in relation to {PHQ}-9 and
{GAD}-7. \emph{Health and Quality of Life Outcomes}, \emph{19}, 260.
\url{https://doi.org/10.1186/s12955-021-01882-x}

\bibitem[\citeproctext]{ref-steger_measuring_2012}
Steger, M. F., Dik, B. J., \& Duffy, R. D. (2012). Measuring meaningful
work: {The} {Work} and {Meaning} {Inventory} ({WAMI}). \emph{Journal of
Career Assessment}, \emph{20}(3), 322--337.
\url{https://doi.org/10.1177/1069072711436160}

\bibitem[\citeproctext]{ref-tennant_warwick-edinburgh_2007}
Tennant, R., Hiller, L., Fishwick, R., Platt, S., Joseph, S., Weich, S.,
Parkinson, J., Secker, J., \& Stewart-Brown, S. (2007). The
{Warwick}-{Edinburgh} {Mental} {Well}-being {Scale} ({WEMWBS}):
Development and {UK} validation. \emph{Health and Quality of Life
Outcomes}, \emph{5}, 63. \url{https://doi.org/10.1186/1477-7525-5-63}

\bibitem[\citeproctext]{ref-tollefson_top_2021}
Tollefson, J. (2021). Top climate scientists are sceptical that nations
will rein in global warming. \emph{Nature}, \emph{599}(7883), 22--24.
\url{https://doi.org/10.1038/d41586-021-02990-w}

\end{CSLReferences}

\appendix

\section{This Section Is an Appendix}\label{apx-a}

\section{Another Appendix}\label{apx-b}






\end{document}